\begin{ledgroupsized}[r]{120mm}
\footnotesize
\pstart
\noindent\textbf{Überlieferung:}
\pend
\end{ledgroupsized}
\begin{ledgroupsized}[r]{114mm}
\footnotesize
\pstart
\parindent -6mm
\makebox[6mm][l]{\textit{L}}%
Konzept:
LH~XXXVII~3~Bl.~81. 
Ein Blatt (ca.~19~x~28~cm.);
linker und unterer Rand beschnitten.
Eine Seite auf Bl.~81~r\textsuperscript{o};
Bl.~81~v\textsuperscript{o} leer.
\pend
\end{ledgroupsized}
\selectlanguage{latin}
\vspace{8mm}
\pstart
\centering
\href{\Scan{LH_37_03_081r}}%
{\lbrack81~r\textsuperscript{o}\rbrack}
\hfill
\edtext{%
De vectibus\indexS{vectes conjugati} conjugatis et figuris\indexS{figurae dentium} dentium in \indexS{rota}rotis
\hfill
\phantom{\lbrack81~r\textsuperscript{o}\rbrack}
\linebreak
resumtae
\edtext{priores cogitationes}{%
\lemma{priores cogitationes}\Cfootnote{Siehe die Erläuterung zu \refnote{FnC:37_03_075-076_resumtae}.}}
\textls{scheda 2\textsuperscript{da}.} 25 Octobr. 1679.%
}{\applabel{FnB:37_03_081_scheda-2da}%
\lemma{}%
\Bfootnote{De vectibus \bracks{...} 1679. \textit{erg.~L}}}
\pend
\nobreak
\vspace{0.5\baselineskip}
\pstart
\noindent
\edlabel{37_03_81_1a}Longe secus est\edlabel{37_03_81_1b}
si concipiamus 
\edtext{dentem\indexS{dens inflexus} inflexum \textit{Acde}}{%
\lemma{dentem}%
\Bfootnote{%
\textit{(1)}~\textit{Acde} %
\textit{(2)}~curvilineum %
\textit{(3)}~flexum
\textit{(4)}~inflexum \textit{Acde}~\textit{L}}} 
esse ducentem, rectum vero \textit{Bf} esse ducendum\lbrack,\rbrack\ quia enim motus jam existit in dente \textit{Acde}, imprimitur statim denti \textit{Bf}, qui eum recipit a linea ad ipsum perpendiculari \textit{\ind{3}d\ind{3}e}. Illud tamen considerandum est, non tam rationem habendam an corpus impingens sit perpendiculare, sed potius an ictus sit perpendicularis, quod hic non fit.
\pend
\begin{edfigure}%
\includegraphics[scale=1]{%
gesamttex/edit_VIII-4/images/LH_37_03_081_d1.pdf}
\edcaption{\edtext{\lbrack\textit{Fig.~1}\rbrack}{%
\lemma{\lbrack\textit{Fig.~1}\rbrack}%
\Cfootnote[nonum]{Zwei gestrichene Entwürfe zum Diagramm werden nicht wiedergegeben.}}}
\end{edfigure}
\pstart
Ad frictionem\indexS{frictio} excludendam possemus duas has considerationes conjungere. Nempe 
\edtext{si retenta}{%
\lemma{si}%
\Bfootnote{%
\textit{(1)}~efficeremus curvam quidem %
\textit{(2)}~retento %
\textit{(3)}~retenta~\textit{L}}}
quidem dentis \textit{AC} 
\edtext{curva \textit{cde} ut ante (scilicet sic, ut semper durante motu in puncto concursus tangens curvae dirigatur in punctum \textit{B})}{%
\lemma{curva \textit{cde}}%
\Bfootnote{%
\textit{(1)}~tam, %
\textit{(2)}~tanta (scilicet sic, \lbrack...\rbrack\ dirigatur
\textit{(a)}~) %
\textit{(b)}~in punctum \textit{B}~\textit{L}}} 
sed denti \textit{Bf} detur flexus talis, ut semper tangens\indexS{tangens curvae} curvae \textit{cde} sit ad eum 
\edtext{perpendicularis, verum}{%
\lemma{perpendicularis,}%
\Bfootnote{%
\textit{(1)}~sed %
\textit{(2)}~verum~\textit{L}}}
jam video non opus esse flexu, sed rem esse praestitam jam tum, nempe ipsi \textit{Bf} in \textit{f} praefigatur portio arcus centro \textit{B}, radio \textit{Bf} descripti, in eum semper impinget dens \textit{Acde} perpendiculariter adeoque hoc modo quam minimam possibilem habebit frictionem\indexS{frictio}, modo 
\edtext{acies hujus arcus}{%
\lemma{acies}%
\Bfootnote{%
\textit{(1)}~hujus parti %
\textit{(2)}~hujus arcus~\textit{L}}}
sit probe elaborata et durata.
\pend
\begin{edfigure}%
\includegraphics[scale=1]{%
gesamttex/edit_VIII-4/images/LH_37_03_081_d2.pdf}
\edcaption{\lbrack\textit{Fig.~2}\rbrack}
\end{edfigure}
\pstart
Denique quo perfectius omnia fiant, et ut ictus\indexS{ictus} ipse quam maxime sit perpendicularis ad radium rotae, adhibeantur dentes quam plurimi; credo ergo si dentium 
\edtext{figura in rota ducen\bracks{da}}{%
\lemma{figura}%
\Bfootnote{%
\textit{(1)}~sit %
\textit{(2)}~in rota %
\textit{(a)}~ducenda %
\textit{(b)}~\bars{ducente \textit{ändert Hrsg.}} sit~\textit{L}}}
curvae \textit{cde}, in ducente autem acies\indexS{acies dentis} dentis 
\edtext{probe elaborata,}{%
\lemma{}%
\Bfootnote{probe \bars{probe \textit{streicht Hrsg.}} elaborata,~\textit{L}}} 
et dentes sint quantum commode 
\edtext{sine motu rupturae}{\lemma{}\Bfootnote{sine motu rupturae \textit{erg.~L}}}
licet 
\edtext{multiplicati et exigui, rotam meliorem}{%
\lemma{multiplicati}%
\Bfootnote{%
\textit{(1)}~, eam
\textit{(2)}~seu
\textit{(3)}~et exigui, rotam
\textit{(a)}~melius %
\textit{(b)}~meliorem~\textit{L}}}
haberi non posse.
\pend
\pstart
Quantum ad rotas\indexS{rotae molendinariae} molendinarias \lbrack(\rbrack Kamrath\indexS{Kamrath} und getriebe\indexS{getriebe}\lbrack)\rbrack\ videtur utile esse, ut striae quaedam in solida columna 
lanternae loco incidantur 
\edtext{secundum lineas aptas,}{%
\lemma{secundum}%
\Bfootnote{%
\textit{(1)}~lineam aptam %
\textit{(2)}~lineas aptas,~\textit{L}}}
ut dentes rotae pectinatae commode in ea ascendere vel descendere, et simul ducere possint. Dentium etiam acies obliquitati striarum adaptata erit, puto melius fore strias in ducente quam in ducendo. 
\pend
\pstart
\selectlanguage{ngerman}%
Es weren nichts als Stucke von schnecken linien, so aber nicht weiter ausgefuhret, als vor einen Zahn, ein wenig zu treiben. Der zähne aber sind desto mehr, wie auch der schnecken\indexS{schnecken linien} 
linien, nach menge der schnecken linien und Zähne kan man die proportiones geben. 
\pend
\begin{edfigure}[.5em][1em]
\includegraphics[scale=1]{%
gesamttex/edit_VIII-4/images/LH_37_03_081_d3.pdf}
\edcaption{\lbrack\textit{Fig.~3}\rbrack}
\end{edfigure}
